Recall that a distribution $\mu$ is $X$-uniform if $\Pr_{g \sim
\mu}[g(x)=y]=\Pr_{g \in G}[g(x)=y]$ for all $x,y \in X$. We say a set $S$
\emph{supports $X$-uniformity} if there exists a distribution supported on
$S$ which is $X$-uniform. We first establish that this is a purely geometric
property of $S$.

Let $R_X$ be the representation of the action of $G$ on $X$, that
is, \[
R_X(g)_{x,y}
= 1_{g(x)=y}\,.
\]
Let $U=R_X(U_G) = \E_{g \in G} [R_X(g)]$ denote the matrix
which corresponds to the action on $X$ of the uniform distribution over
$G$. We consider these matrices as points in $\R^{d}$ for $d=|X|^2$.

\begin{proposition}\label{prop:uniform-by-convex-hull}
A set $S \subset G$ supports $X$-uniformity iff the convex
hull of the matrices $\{R_X(g): g \in S\}$ contains the matrix $U$.
\end{proposition}
\begin{proof}
A point in the convex hull is given by $M = \sum_{g \in S} \mu(g) \cdot
R_X(g)$ where $\mu(g) \ge 0$ and $\sum_{g \in S} \mu(g)=1$. Thus, each
point in the convex hull corresponds to a distribution $\mu$ over $S$,
and vice versa. Note that $M_{x,y}=\Pr_{g \sim \mu}[g(x)=y]$, hence an
$X$-uniform distribution corresponds to the matrix $U$.
\end{proof}
